% =====================================================================
% P5 merge snippet — M08: derivations of the balance laws used in this paper
% Source: manuscript_v2, Appendix B (app:balance_derivations): momentum, symmetry of stress,
% mechanical power, internal energy, and the exact Eulerian flux.
% Adapted: this paper's notation (\theta, \bu, \dot u) and equation labels; the paper states
% the balances (eq:mom) but does not derive them.
% NEW LABELS: app:balance, eq:bal-mom, eq:bal-sym, eq:bal-power, eq:bal-energy
% Suggested placement: new appendix before \cref{app:disp}, since \cref{app:disp} uses the
% field equations that these derivations produce.
% =====================================================================

\section{Derivations of the balance laws}\label{app:balance}

The balances in \cref{sec:foundations} are used throughout; this appendix records where each comes from
and what it assumes, so that a reader can see which statements are identities of continuum mechanics and
which already contain a constitutive choice. Every derivation below is algebra on the stated assumptions;
none uses a heat law.

\subsection*{Linear momentum}\label{eq:bal-mom}
Take an arbitrary material volume $V$ with surface $\partial V$ and outward normal $n_i$. Cauchy's
postulate gives surface traction $t_i=\sigma_{ji}n_j$. Requiring that the rate of change of linear momentum
equal the resultant force for \emph{every} $V$,
$\frac{\dd}{\dd t}\int_V\rho\dot u_i\,\dd V=\int_{\partial V}\sigma_{ji}n_j\,\dd S+\int_V\rho f_i\,\dd V$,
and applying the divergence theorem together with the transport theorem gives
$\rho\ddot u_i=\sigma_{ji,j}+\rho f_i$ after the volume is shrunk to a point. Assumptions: a Cauchy
stress exists (no couple stresses), motion is described by a single displacement field, and the body is
non-polar. These are exactly the assumptions recorded for \eqref{eq:mom}.

\subsection*{Symmetry of the stress}\label{eq:bal-sym}
Applying the same argument to the moment of momentum about a point, with no body couples and no couple
stresses, leaves $\int_V e_{ijk}\sigma_{jk}\,\dd V=0$ for every $V$, hence $\sigma_{ij}=\sigma_{ji}$. The
conclusion is a consequence of the assumptions, not an independent law: in a micropolar or couple-stress
continuum the stress is not symmetric and the derivation must be redone with the couple-stress tensor.

\subsection*{Mechanical power}\label{eq:bal-power}
Contracting \eqref{eq:bal-mom} with the velocity and using the symmetry of the stress,
\begin{equation}
\rho\frac{\dd}{\dd t}\Bigl(\tfrac12\dot u_i\dot u_i\Bigr)=\bigl(\sigma_{ij}\dot u_i\bigr)_{,j}
-\sigma_{ij}\dot\varepsilon_{ij}+\rho f_i\dot u_i ,
\end{equation}
so the stress power $\sigma_{ij}\dot\varepsilon_{ij}$ is the only term that exchanges mechanical energy with
the internal state. This is where the strain, and not the displacement gradient, enters --- the rotation
does no work.

\subsection*{Internal energy}\label{eq:bal-energy}
Postulating that the internal energy of a material volume changes by the heat supplied plus the stress power
delivered, with heat flux $q_i$ and supply $r$,
\begin{equation}
\rho\dot e=\sigma_{ij}\dot\varepsilon_{ij}-q_{i,i}+\rho r ,
\end{equation}
which is the second line of \eqref{eq:mom}. This is an axiom of the theory: it defines what $e$ means, and
a generalized theory that changes the flux $q_i$ or adds a state variable must restate it.

\subsection*{The exact Eulerian flux}
Adding the power balance and the internal-energy balance gives the material-derivative form
$\rho\frac{\dd}{\dd t}\bigl(\tfrac12\dot u_i\dot u_i+e\bigr)+\bigl(q_j-\sigma_{ji}\dot u_i\bigr)_{,j}
=\rho f_i\dot u_i+\rho r$. Converting the material derivative with mass conservation yields
\eqref{eq:exactflux}. The step that is often skipped in the literature is the last one: the material form
is not the fixed-volume form, and the difference is the advection term $\mathcal E_{\mathrm{tot}}\dot u_j$.
Omitting it is legitimate only as a first-order reduction under the small-motion assumptions stated with
\eqref{eq:exactflux-balance}.
